\documentclass[aspectratio=169,10pt]{beamer}
% =====================================================================
%  Le Collé — thème de présentation (même palette et mêmes notations que la fiche)
% =====================================================================
\usepackage[utf8]{inputenc}
\usepackage[T1]{fontenc}
\usepackage[sc]{mathpazo}
\usepackage[scaled=0.92]{helvet}
\usefonttheme{professionalfonts}
\usepackage{amsmath,amssymb,mathrsfs,stmaryrd}
\usepackage{microtype}
\usepackage{booktabs,array,multicol}
\usepackage[most]{tcolorbox}
\usepackage{tikz,pgfplots}
\usetikzlibrary{arrows.meta,calc,positioning,decorations.pathreplacing,shapes.misc}
\pgfplotsset{compat=1.17,trig format plots=rad}
\usepackage{tkz-tab}

% ---------- palette (identique à la fiche) ----------
\definecolor{Ink}{HTML}{1C2533}\definecolor{Bleu}{HTML}{1F3A5F}\definecolor{BleuMoyen}{HTML}{4A6A8F}
\definecolor{BleuClair}{HTML}{9DB2C9}\definecolor{Fond}{HTML}{F2F4F7}\definecolor{Ocre}{HTML}{B07A26}
\definecolor{Brique}{HTML}{A23B2A}\definecolor{Gris}{HTML}{5E6672}\definecolor{Filet}{HTML}{C5CCD6}
\definecolor{Vert}{HTML}{2E6B4F}

% ---------- notations (identiques à la fiche) ----------
\newcommand{\K}{\mathbb{K}}\newcommand{\R}{\mathbb{R}}\newcommand{\C}{\mathbb{C}}\newcommand{\Z}{\mathbb{Z}}
\newcommand{\Mn}[1]{\mathcal{M}_{n}(#1)}\newcommand{\rg}{\operatorname{rg}}\newcommand{\vect}{\operatorname{Vect}}
\newcommand{\Mat}{\operatorname{Mat}}\newcommand{\Id}{\operatorname{Id}}\newcommand{\com}{\operatorname{com}}
\newcommand{\vv}[1]{\overrightarrow{#1}}\newcommand{\ve}[1]{\vec{#1}}
\newcommand{\dt}[1]{\left|\begin{matrix}#1\end{matrix}\right|}
\DeclareRobustCommand{\HP}{{\sffamily\bfseries\scriptsize\textcolor{Brique}{[HP]}}}

% ---------- réglages beamer ----------
\setbeamertemplate{navigation symbols}{}
\setbeamersize{text margin left=0.85cm,text margin right=0.85cm}
\setbeamercolor{normal text}{fg=Ink,bg=white}
\setbeamercolor{structure}{fg=Bleu}
\setbeamercolor{alerted text}{fg=Brique}
\setbeamerfont{alerted text}{series=\bfseries}
\setbeamertemplate{itemize item}{\raisebox{0.15em}{\textcolor{Ocre}{\rule{0.38em}{0.38em}}}}
\setbeamertemplate{itemize subitem}{\textcolor{BleuMoyen}{\scriptsize\textbullet}}
\setbeamertemplate{enumerate item}{\textcolor{Bleu}{\sffamily\bfseries\insertenumlabel.}}
\setlength{\leftmargini}{1.2em}
\setbeamertemplate{blocks}[default]

% numéro de la partie courante, rappelé dans le titre et le pied de page (renvoi à la fiche)
\def\fichesec{}\def\chapnom{}
\newcommand{\chapitre}[1]{\def\chapnom{#1}}

\setbeamertemplate{frametitle}{%
  \vskip0.28cm
  \ifx\fichesec\empty\else
    \tikz[baseline=(b.base)]\node[fill=Ocre,text=white,font=\sffamily\bfseries\scriptsize,inner sep=2.5pt,rounded corners=1pt](b){\fichesec};\hspace{0.45em}%
  \fi
  {\color{Bleu}\sffamily\bfseries\Large\insertframetitle}\par
  \vskip0.06cm{\color{Filet}\rule{\linewidth}{0.5pt}}\vskip-0.2cm}

\setbeamertemplate{footline}{%
  \hbox to\paperwidth{%
    \hspace*{0.85cm}{\scriptsize\sffamily\color{Gris}Le Collé\ \ \textcolor{Filet}{|}\ \ \chapnom}\hfill
    {\scriptsize\sffamily\color{Gris}\ifx\fichesec\empty\else fiche \textbf{\color{Ocre}§\fichesec}\quad\fi\insertframenumber}\hspace*{0.85cm}}%
  \vskip0.18cm
  \hbox{\begin{tikzpicture}
    \fill[Fond] (0,0) rectangle (\paperwidth,2.2pt);
    \pgfmathsetmacro\prog{min(\insertframenumber/max(\inserttotalframenumber,1),1)}
    \fill[Ocre] (0,0) rectangle ({\prog*\paperwidth},2.2pt);
  \end{tikzpicture}}}

% ---------- encadrés : même code couleur que la fiche ----------
\tcbset{sb/.style={enhanced,arc=1pt,boxrule=0pt,left=6pt,right=6pt,top=3pt,bottom=3pt,before skip=0.14cm,after skip=0.14cm,
  fontupper=\normalsize\raggedright}}
\newcommand{\etiq}[2]{{\sffamily\bfseries\scriptsize\color{#1}\MakeUppercase{#2}}\par\vspace{2pt}}
\newtcolorbox{DEF}[1][Définition]{sb,colback=white,borderline west={2pt}{0pt}{Bleu},before upper={\etiq{Bleu}{#1}}}
\newtcolorbox{PROP}[1][Propriété]{sb,colback=Fond,borderline west={3pt}{0pt}{Bleu},before upper={\etiq{Bleu}{#1}}}
\newtcolorbox{FORM}[1][Formule]{sb,colback=white,colframe=BleuMoyen,boxrule=0.6pt,before upper={\etiq{BleuMoyen}{#1}}}
\newtcolorbox{METH}[1]{sb,colback=white,borderline north={1pt}{0pt}{Bleu},borderline south={1pt}{0pt}{Bleu},
  before upper={{\sffamily\bfseries\scriptsize\colorbox{Bleu}{\color{white}MÉTHODE}\ \color{Bleu}#1}\par\vspace{3pt}}}
\newtcolorbox{PIEGE}[1][Piège]{sb,colback=Brique!6,borderline west={2pt}{0pt}{Brique},before upper={\etiq{Brique}{#1}}}
\newtcolorbox{ASTU}[1][Astuce]{sb,colback=Ocre!8,borderline west={2pt}{0pt}{Ocre},before upper={\etiq{Ocre}{#1}}}
\newtcolorbox{HPBOX}{sb,colback=white,borderline west={2pt}{0pt}{Brique},before upper={\etiq{Brique}{Hors programme}}}
\newtcolorbox{ATOI}[1][À toi]{sb,colback=Vert!6,borderline west={2pt}{0pt}{Vert},before upper={\etiq{Vert}{#1}}}
\newtcolorbox{RETENIR}{sb,colback=Bleu!5,colframe=Bleu,boxrule=0.8pt,left=8pt,
  before upper={\etiq{Bleu}{À retenir}}}
\newtcolorbox{EX}[1][Exemple]{sb,colback=white,colframe=Filet,boxrule=0.5pt,before upper={\etiq{Gris}{#1}}}

% complément pédagogique absent du livre : étiquette pointillée
\DeclareRobustCommand{\ajout}{\tikz[baseline=(a.base)]\node[draw=Gris,dashed,text=Gris,font=\sffamily\scriptsize,inner sep=2pt,rounded corners=1pt](a){complément hors livre};}
\DeclareRobustCommand{\livre}[1]{{\sffamily\scriptsize\color{Gris}[livre #1]}}
\newcommand{\rep}[1]{\textcolor{Vert}{#1}}% réponse attendue de l'élève

% ---------- styles de figures (identiques au livre) ----------
\pgfplotsset{courbe/.style={axis lines=middle,axis equal image,
  axis line style={Gris,-{Stealth[length=4pt]}},
  xlabel={$x$},ylabel={$y$},ticklabel style={font=\scriptsize\sffamily,color=Gris},
  every axis x label/.style={at={(ticklabel* cs:1)},anchor=west,font=\small},
  every axis y label/.style={at={(ticklabel* cs:1)},anchor=south,font=\small},
  samples=200,smooth}}
\pgfplotsset{mini/.style={axis lines=middle,axis line style={Filet},xtick=\empty,ytick=\empty,width=3.9cm,height=3.9cm,
  xmin=-1.15,xmax=1.15,ymin=-1.15,ymax=1.15,title style={font=\footnotesize\sffamily,yshift=-0.1cm},samples=100}}
\tikzset{arcetude/.style={Bleu,line width=1.25pt},
         reste/.style={BleuClair,line width=1pt},
         tgte/.style={Ocre,line width=1pt},
         asy/.style={Gris,dashed,line width=0.7pt},
         pt/.style={circle,fill=Ink,inner sep=1.3pt},
         >={Stealth[length=5pt]}}

% ---------- diapositives de partie ----------
% \partie{num}{titre}{question guide}{liste des parties du chapitre via \rm}
\newcommand{\rmitem}[2]{\def\tmp{#1}%
  \ifx\tmp\fichesec{\color{Bleu}\bfseries #1\quad #2}\else{\color{Gris!70}#1\quad #2}\fi\par\vspace{2pt}}
\newcommand{\partie}[4]{%
  \def\fichesec{#1}\section{#2}%
  \begin{frame}[plain]
  \begin{tikzpicture}[remember picture,overlay]
    \fill[Bleu] (current page.north west) rectangle ([xshift=6.3cm]current page.south west);
    \node[anchor=north west,text=white,font=\sffamily\bfseries\fontsize{40}{40}\selectfont] at ([xshift=0.8cm,yshift=-1.2cm]current page.north west) {#1};
    \node[anchor=north west,text width=4.8cm,align=left,text=white,font=\sffamily\bfseries\Large] at ([xshift=0.8cm,yshift=-2.9cm]current page.north west) {\hyphenpenalty=10000\exhyphenpenalty=10000 #2\par};
    \node[anchor=south west,text width=4.8cm,text=white!75!Bleu,font=\sffamily\scriptsize] at ([xshift=0.8cm,yshift=0.6cm]current page.south west) {\chapnom\\fiche \textbf{§#1}};
  \end{tikzpicture}
  \hspace*{5.5cm}\begin{minipage}{8.4cm}
  {\sffamily\scriptsize\color{Ocre}\bfseries LA QUESTION}\par\vspace{2pt}
  {\large #3}\par\vspace{0.7cm}
  {\sffamily\scriptsize\color{Gris}\bfseries DANS CE CHAPITRE}\par\vspace{4pt}
  {\sffamily\small #4}
  \end{minipage}
  \end{frame}}



% ---------- exercices en séance : énoncé sur une diapositive, correction sur la suivante ----------
\newcommand{\badgeexo}[1]{\tikz[baseline=(e.base)]\node[fill=Vert,text=white,font=\sffamily\bfseries\scriptsize,inner sep=2.5pt,rounded corners=1pt](e){EXERCICE #1};\hspace{0.4em}}
\newcommand{\badgecor}[1]{\tikz[baseline=(e.base)]\node[draw=Vert,line width=0.8pt,text=Vert,font=\sffamily\bfseries\scriptsize,inner sep=2.2pt,rounded corners=1pt](e){CORRECTION #1};\hspace{0.4em}}
\newtcolorbox{ENONCE}[1][]{sb,colback=Vert!6,borderline west={3pt}{0pt}{Vert},left=9pt,top=5pt,bottom=5pt,
  before upper={{\sffamily\bfseries\scriptsize\color{Vert}À TOI\ifx\relax#1\relax\else\ \ ·\ \ \MakeUppercase{#1}\fi}\par\vspace{4pt}}}
\newtcolorbox{INDICE}{sb,colback=white,colframe=Vert!60,boxrule=0.5pt,borderline={0.5pt}{0pt}{Vert!60,dashed},
  before upper={\etiq{Vert}{Indice (si besoin)}}}
\newtcolorbox{RAIS}[1][Le raisonnement]{sb,colback=Vert!5,borderline west={2pt}{0pt}{Vert},before upper={\etiq{Vert}{#1}}}
\newtcolorbox{OUTIL}{sb,colback=Fond,borderline west={2pt}{0pt}{Bleu},before upper={\etiq{Bleu}{Propriété ou méthode utilisée}}}
\newcommand{\etape}[1]{{\sffamily\bfseries\footnotesize\color{Vert}#1}\ }
\newcommand{\objectif}[1]{\par\vspace{2pt}{\footnotesize\sffamily\color{Gris}Objectif : #1}}
